% =====================================================================
% Fichas de evidencia del encargo E1 (sección 4.2 de la guía).
% Una ficha por inciso. Los números deben coincidir con results.json
% regenerado desde cero (se sincronizarán al cerrar el encargo).
% =====================================================================
\section{Fichas de evidencia}

% ---------------------------------------------------------------------
\begin{recuadro}[title=Ficha de evidencia — E1.1: tres métodos de solución de la ecuación de Riccati]
\begin{tabularx}{\linewidth}{@{}>{\bfseries}p{3.1cm} >{\raggedright\arraybackslash}X@{}}
Afirmación &
Los métodos (a) vectores propios de la Hamiltoniana, (b) Schur real ordenado y (c) Newton--Kleinman
resuelven la ecuación de Riccati del diseño base $(Q_0,R_0)$ con una diferencia máxima entre ellos de
$2.5\times10^{-12}$. En el banco de seis problemas, (a) y (c) quedan por debajo de $10^{-10}$ frente a
\texttt{icare} en todos los casos; (b) supera esa tolerancia solo en el modelo de 5 estados con armadura
($7.2\times10^{-9}$), cuya $P$ es casi singular ($\lambda_{\min}(P)=7.4\times10^{-14}$). \\[3pt]
Evidencia &
Tablas de \texttt{e1\_1\_metodos.m}; figura \texttt{figuras/E1\_1\_errores.png};
datos en \texttt{resultados/E1\_1.mat}. Implementación en \texttt{src/riccati\_hamiltoniano.m},
\texttt{src/riccati\_schur.m}, \texttt{src/riccati\_newton.m}. \\[3pt]
Comprobación independiente &
(1) Biblioteca: \texttt{icare}; diseño base: (a) $9.7\times10^{-15}$, (b) $2.4\times10^{-12}$,
(c) $1.0\times10^{-13}$.
(2) Kit: $P_0$ de \texttt{diseno\_base\_ref.json} (\texttt{pruebas\_E1.m}); $K$ de Newton frente a
$K_0$: $5.4\times10^{-14}$.
(3) Respuesta conocida: $\dot x=x+u$, $K=1+\sqrt{1+q}$; error $\le 2.2\times10^{-16}$ para
$q\in[10^{-6},1]$.
(4) Convergencia cuadrática de Newton--Kleinman: residuo
$1.7\times10^{-3}\to3.6\times10^{-5}\to1.2\times10^{-8}\to1.2\times10^{-15}$ (iteraciones 6 a 9). \\[3pt]
Supuestos &
$(A,B)$ estabilizable y $(A,Q^{1/2})$ detectable. Newton requiere una ganancia inicial
estabilizante: Bass con $\beta=20$ en 4 estados; en 5 estados Bass exigiría $\beta>7240$ y no
estabiliza numéricamente, por lo que se usa $[K_0\;0]$. El método (b) se implementó sin escalado
ni simetrización de $P$; con balanceo (como \texttt{icare}) su error debería disminuir, lo que
cambiaría la ordenación de los métodos.
\end{tabularx}
\end{recuadro}

% ---------------------------------------------------------------------
\begin{recuadro}[title=Ficha de evidencia — E1.2: verificación de las soluciones y condicionamiento]
\begin{tabularx}{\linewidth}{@{}>{\bfseries}p{3.1cm} >{\raggedright\arraybackslash}X@{}}
Afirmación &
Un residuo pequeño no implica un error pequeño en $P$: en el oscilador
$A=\left[\begin{smallmatrix}0&1\\-1&0\end{smallmatrix}\right]$, $Q=\varepsilon^2 I$, una $P$ con error relativo
de $1\,\%$ tiene residuo normalizado de solo $7.1\times10^{-9}$ para $\varepsilon=10^{-6}$, porque
$\|\mathcal{L}^{-1}\|_2=7.07\times10^{5}$, con $\mathcal{L}(X)=A_c^\top X+XA_c$. En el péndulo,
$\|\mathcal{L}^{-1}\|_2=86$ y la cota $\|\mathcal{L}^{-1}\|\,\|\mathrm{Res}\|/\|P\|\le 6.9\times10^{-10}$
acota el error real ($\le 2.4\times10^{-12}$). Las tres soluciones del péndulo pasan las cuatro
pruebas: residuo $\le1.4\times10^{-13}$, asimetría $\le1.9\times10^{-13}$,
$\lambda_{\min}(P)=3.4\times10^{-4}>0$ y $\max\operatorname{Re}\lambda(A-BK)=-2.668$. \\[3pt]
Evidencia &
Tablas de \texttt{e1\_2\_verificacion.m}; figura \texttt{figuras/E1\_2\_residuo\_vs\_error.png};
datos en \texttt{resultados/E1\_2.mat}; \texttt{src/operador\_lyapunov.m}, \texttt{src/verificar\_riccati.m}. \\[3pt]
Comprobación independiente &
(1) Solución analítica del oscilador: $p_2=\varepsilon^2/(1+\sqrt{1+\varepsilon^2})$,
$p_3=\sqrt{2p_2+\varepsilon^2}$, $p_1=p_3(1+p_2)$.
(2) Primer orden: $\mathrm{Res}(P^\star+\Delta)\approx\mathcal{L}(\Delta)$ con discrepancia de $0.7\,\%$ en
todas las filas (término cuadrático $\Delta S\Delta$).
(3) Control: en una dirección aleatoria el mismo error da residuo $3.4\times10^{-3}$.
(4) La propia \texttt{icare} pierde precisión frente a la solución exacta: de $1.8\times10^{-14}$
($\varepsilon=10^{-1}$) a $2.1\times10^{-6}$ ($\varepsilon=10^{-6}$). \\[3pt]
Supuestos &
Análisis de primer orden (válido si $\|\Delta S\Delta\|\ll\|\mathcal{L}(\Delta)\|$, aquí $0.7\,\%$).
La dirección ``peor'' es la del menor valor singular de $\mathcal{L}$ (caso pesimista); otra dirección
da un residuo mayor. La normalización del residuo es la de \texttt{residuo\_riccati.m}; otra
normalización cambia los valores, no la tendencia. $\|\mathcal{L}^{-1}\|$ no explica por completo el
error de \texttt{icare}, que crece más rápido; el condicionamiento completo de la ecuación de Riccati
(Kenney, Laub y Wette, 1990) incluye otros términos.
\end{tabularx}
\end{recuadro}

% ---------------------------------------------------------------------
\begin{recuadro}[title=Ficha de evidencia — E1.3: estudio de estrés numérico]
\begin{tabularx}{\linewidth}{@{}>{\bfseries}p{3.1cm} >{\raggedright\arraybackslash}X@{}}
Afirmación &
La precisión la decide el \emph{escalado de los datos}, no el método de subespacio. Sin balanceo,
vectores propios y Schur pierden hasta 12 cifras con $R=10^{-8}$ ($1.2\times10^{-4}$) y 14 con las
velocidades en otras unidades ($s=10^{6}$: $3.0\times10^{-2}$). Balancear $H$ recupera F2
($1.6\times10^{-14}$), pero en F1 deja $3.6\times10^{-9}$: formar $S=BR^{-1}B^\top$ ya mezcla escalas
incompatibles. \texttt{icare} evita ambos problemas: usa el haz extendido (no forma $S$) y escala
estados y entrada; reproducido con código propio, alcanza su precisión (peor caso $9.8\times10^{-14}$).
Con autovalores de $H$ cerca del eje imaginario, el error global sigue $\le2.4\times10^{-12}$, pero el
de la entrada pequeña $P_{11}$ llega a $2.6\times10^{-9}$. \\[3pt]
Evidencia &
\texttt{e1\_3\_estres.m} (tres familias, 39 problemas), \texttt{e1\_3b\_balanceo.m}, \texttt{e1\_3c\_icare.m},
\texttt{e1\_3d\_escalado.m}, \texttt{e1\_3e\_icare\_escalado.m}; figuras \texttt{E1\_3\_*.png}; datos en
\texttt{resultados/E1\_3*.mat}; \texttt{src/riccati\_pencil.m}, \texttt{src/escalar\_riccati.m}. \\[3pt]
Comprobación independiente &
(1) Referencia en precisión extendida: Newton--Kleinman con \texttt{vpa} de 40 dígitos y datos exactos
(\texttt{sym(X,'f')}).
(2) Experimento $2\times2$ \{vectores propios, Schur\} $\times$ \{con, sin balanceo\}: las columnas con el
mismo balanceo coinciden en las 39 filas.
(3) Haz extendido propio sin escalado frente a \texttt{icare(\ldots,'noscaling')}: F3 idénticos fila a fila
(p.\ ej.\ $8.8\times10^{-10}$ ambos), F2 $6.4\times10^{-5}$ frente a $6.3\times10^{-5}$ ($s=10^5$); con $s=10^6$
\texttt{icare} sin escalado no devuelve solución.
(4) El mismo haz con los vectores de escalado de \texttt{icare} (\texttt{info.Sx}, \texttt{info.Sr}):
peor caso F1/F2/F3 $9.8\times10^{-14}$, $8.4\times10^{-15}$, $1.6\times10^{-14}$, frente a \texttt{icare}
$2.6\times10^{-13}$, $7.4\times10^{-15}$, $2.0\times10^{-14}$.
(5) Hipótesis descartadas: simetrizar $P$ y $\operatorname{cond}(U_{11})$ (vale 4.0 en el peor caso).
(6) Newton--Kleinman cuadrático frente a \texttt{vpa}: $e_k/e_{k-1}^2=0.81,\,1.03,\,1.07$. \\[3pt]
Supuestos &
Se reutiliza el escalado que elige \texttt{icare}, no su regla de elección. Los escalados propios
ingenuos empeoran: igualar $\|S\|$ y $\|Q\|$ ignora $A$ ($5.3\times10^{-8}$ en F1), y \texttt{balance(A)}
degenera porque la primera columna de $A$ es nula ($\theta$ no aparece en la dinámica), con
$D=\operatorname{diag}(1,\,6\times10^{-8},\,1,\,1)$. El error en norma de Frobenius oculta el de entradas
pequeñas; por eso se reporta también el de $P_{11}$. Cada familia varía un solo parámetro.
\end{tabularx}
\end{recuadro}

% ---------------------------------------------------------------------
\begin{recuadro}[title=Ficha de evidencia — E1.4: frente de Pareto]
\begin{tabularx}{\linewidth}{@{}>{\bfseries}p{3.1cm} >{\raggedright\arraybackslash}X@{}}
Afirmación &
Los 625 diseños LQR obtenidos al escalar los pesos de brazo y péndulo entre $10^{-3}$ y $10^{3}$ (normalización
de Bryson del diseño base) son todos estables y no dominados, y forman una superficie de Pareto con
$J_\theta/(\pi/4)^2\in[8.96\times10^{-3},\,2.23\times10^{2}]$, $J_\alpha/(\pi/36)^2\in[2.03\times10^{-3},\,0.141]$ y
$J_u/10^2\in[1.56\times10^{-3},\,1.11]$. En la planta no lineal, los 625 tienen éxito en C1, pero en C2 solo
147 (estado real) y 115 (estimador base, cuantización y ruido); 101 cumplen además una fracción saturada
$\le5\,\%$. 48 diseños exitosos con estado real fallan con el sensor. \\[3pt]
Evidencia &
\texttt{e1\_4\_pareto.m}; figuras \texttt{figuras/E1\_4\_pareto3d.png} y \texttt{figuras/E1\_4\_proyecciones.png};
datos en \texttt{resultados/E1\_4.mat} y simulaciones en \texttt{resultados/E1\_4\_sim.mat}. \\[3pt]
Comprobación independiente &
(1) Lyapunov por dos caminos: primal (gramiano, \texttt{sylvester}) y dual (una ecuación por objetivo,
\texttt{lyap\_kron}): diferencia $\le1.3\times10^{-9}$.
(2) Respuesta conocida: la suma ponderada de los objetivos coincide con $\operatorname{tr}(PX_0)$ del LQR
($\le1.9\times10^{-10}$), y el filtro de Pareto declara no dominados 625 de 625, como predice la teoría.
(3) Simulación lineal con \texttt{ode45} en 10 diseños: diferencia $\le2.3\times10^{-11}$.
(4) \texttt{icare} frente a la Hamiltoniana balanceada: $\le8.9\times10^{-11}$.
(5) El diseño base reproduce $K_0$ y los valores del kit: 9.12~V y 1.237~rad (estado real), 63.82~V y
1.345~rad (con estimador).
(6) El esfuerzo de Lyapunov ordena el esfuerzo no lineal: correlación de rangos de Spearman 0.960 entre
$J_u$ y $\max|u_c|$ en C1. \\[3pt]
Supuestos &
Covarianza inicial declarada $X_0=\operatorname{diag}(0.1^2,\,0.05^2,\,0.5^2,\,0.5^2)$ (unidades SI): otra $X_0$
cambia la forma del frente, no la no dominancia. Los objetivos esperados son lineales y no ven la
saturación; por eso la factibilidad se juzga en la planta no lineal. C2 con sensor se simuló con una sola
semilla de ruido: los 64 diseños que cambian de resultado entre estado real y sensor están cerca del límite
y podrían cambiar con otra semilla. La malla varía dos pesos con $R$ fija; escalar $Q$ y $R$ a la vez no
cambia $K$.
\end{tabularx}
\end{recuadro}

% ---------------------------------------------------------------------
\begin{recuadro}[title=Ficha de evidencia — E1.5: optimización de los pesos con restricciones]
\begin{tabularx}{\linewidth}{@{}>{\bfseries}p{3.1cm} >{\raggedright\arraybackslash}X@{}}
Afirmación &
Con las restricciones del grupo (decaimiento $\sigma\ge0.458$~s$^{-1}$, amortiguamiento $\zeta\ge0.5$, C2 con
$|\theta|\le\pi/2$ y saturación $\le5\,\%$, éxito en C1--C3b), ningún LQR mejora de forma significativa el
costo oficial promedio del diseño base: el óptimo por promedio de 5 semillas ($r_\theta=1.58$, $r_\alpha=1$)
da $J_{rel}=0.9994\pm0.0225$ en 20 semillas independientes ($t=0.12$). Hay un compromiso significativo:
penalizar más el brazo mejora C2 ($+1.5\,\%$, $t=45.8$) y empeora C1 ($-1.2\,\%$, $t=-4.6$). El diseño base
es factible ($\sigma=2.668$, $\zeta=0.742$, saturación 1.9\,\%, $\max|\theta|=1.345$~rad en C2) y se adopta
como diseño final. \\[3pt]
Evidencia &
\texttt{e1\_5\_restricciones.m}, \texttt{e1\_5b\_saa.m}, \texttt{e1\_5c\_sensibilidad.m}, \texttt{src/diseno\_final.m};
figuras \texttt{figuras/E1\_5\_*.png}; datos en \texttt{resultados/E1\_5.mat} y \texttt{E1\_5b.mat}. \\[3pt]
Comprobación independiente &
(1) Validación con semillas 101--120, no usadas para elegir, en comparación pareada con el base (prueba
$t$, 19 grados de libertad).
(2) Dos búsquedas \texttt{fminsearch} desde puntos distintos llegan al mismo diseño ($J_{rel}=0.9625$ y
$0.9614$ con la semilla 1).
(3) Piso de ruido del costo oficial: perturbar $K_0$ en $10^{-15}$ cambia $J$ en 0.55\,\% (C1), 0.01\,\% (C2),
3.1\,\% (C3) y 0.09\,\% (C3b) con cuantización, y en $\sim10^{-16}$ sin ella; explica el $J_{rel}=0.9991$ del base
recalculado con \texttt{icare}.
(4) $\sigma\ge0.458$ es coherente con la simulación: tiempo de establecimiento en C1 $\le0.132$~s.
(5) La optimización con una sola semilla sobreestimó la mejora ($0.961$ frente a $0.9945\pm0.019$). \\[3pt]
Supuestos &
$J_{rel}$ pondera igual los cuatro escenarios: dar más peso a C2 y C3b favorecería un diseño más agresivo
(el óptimo de la semilla 1 mejora C3b en $5.2\,\%$, $t=8.3$, pero empeora C1 en $8.2\,\%$). Búsqueda en LQR con
$Q$ diagonal; los pesos de las velocidades no ayudaron (\texttt{fminsearch} los llevó al mínimo). $\sigma_{\min}$
proviene de una envolvente exponencial aproximada. La prueba $t$ supone diferencias aproximadamente normales.
\end{tabularx}
\end{recuadro}

% ---------------------------------------------------------------------
\begin{recuadro}[title=Ficha de evidencia — E1.5 (bonificación): síntesis por regiones de polos con LMI]
\begin{tabularx}{\linewidth}{@{}>{\bfseries}p{3.1cm} >{\raggedright\arraybackslash}X@{}}
Afirmación &
Imponer la región del grupo ($\sigma\ge0.458$, $\zeta\ge0.5$) por LMI con Lyapunov común cuesta $+2.5\,\%$ en
$J_0$ aunque $K_0$ ya está en la región: el método es conservador (logra $\sigma=2.939$, $\zeta=0.885$). Frente
al ajuste de los dos pesos de posición de $Q$, el ajuste de $Q$ es más barato al exigir decaimiento
($\sigma\ge4$: $1.18$ frente a $1.57$; $\sigma\ge6$: $1.46$ frente a $2.79$, en $J_0/J_{0,base}$) y la LMI lo es al
exigir amortiguamiento ($\zeta\ge0.85$: $1.18$ frente a $1.48$; $\zeta\ge0.9$: $1.25$ frente a $1.98$). Una región
alcanzable en el modelo lineal no garantiza la planta saturada: con $\sigma\ge4$ el diseño cumple C1 pero
falla C2 incluso con estado real ($99.9\,\%$ del tiempo saturado). \\[3pt]
Evidencia &
\texttt{e1\_5d\_lmi.m}, \texttt{src/lqr\_region\_lmi.m} (LMI Lab); figura \texttt{figuras/E1\_5\_lmi.png};
datos en \texttt{resultados/E1\_5\_lmi.mat}. \\[3pt]
Comprobación independiente &
(1) Respuesta conocida: sin región, la LMI reproduce el LQR ($\|K-K_0\|/\|K_0\|=1.3\times10^{-4}$, costo
$1.000000\,J_{0,base}$).
(2) Los polos de cada diseño, calculados con \texttt{eig}, cumplen la región pedida en los 16 casos.
(3) La cota $\operatorname{tr}(Z)$ queda por encima del costo real $J_0$ (Lyapunov propio) en todos los casos.
(4) Validación no lineal con semillas 101--120: el diseño LMI de la región del grupo cumple R3--R4 en
20 de 20, pero no mejora a $K_0$ ($J_{rel}=0.9982\pm0.0216$, $t=0.37$): mejora C3b ($+7.6\,\%$) y empeora C1
($-6.6\,\%$) y C2 ($-3.2\,\%$). \\[3pt]
Supuestos &
Una sola matriz de Lyapunov para el costo y la región (formulación de Chilali y Gahinet): la cota es
conservadora. La comparación con el ajuste de $Q$ usa solo los pesos de posición de la malla de E1.4; con
pesos en las velocidades el LQR podría amortiguar más barato. Con 5 semillas (1--5) el diseño LMI parecía
$2.3\,\%$ mejor ($0.977\pm0.009$); no se confirmó con 20 semillas independientes. Precisión del solucionador
LMI: $\sim10^{-4}$ en $K$.
\end{tabularx}
\end{recuadro}

% ---------------------------------------------------------------------
\begin{recuadro}[title=Ficha de evidencia — E1.6: región de atracción del regulador saturado]
\begin{tabularx}{\linewidth}{@{}>{\bfseries}p{3.1cm} >{\raggedright\arraybackslash}X@{}}
Afirmación &
\emph{Certificado} (modelo lineal sin saturación): $\rho^\star=V_{\max}^2/(KP^{-1}K^\top)=1.635\times10^{-2}$; el
elipsoide $\{x^\top Px\le\rho^\star\}$ llega a $0.0377$~rad sobre el eje $\alpha$ y a $0.144$~rad sobre $\theta$.
\emph{Estimado} (planta no lineal saturada, lazo digital): el péndulo se recupera desde $\alpha_0=0.430$~rad
(estado real) y cumple además $|\theta|\le\pi/2$ hasta $0.388$~rad; con estimador, cuantización y ruido,
hasta $0.378$~rad ($-2.4\,\%$). El certificado subestima el radio entre 8.5 y 32 veces (mediana 12.2 en 200
direcciones). En la caja declarada se recupera el $44.8\,\%$ de 20\,000 condiciones iniciales ($34.6\,\%$ sin
exceder $|\theta|\le\pi/2$). C2 ($\alpha_0=0.35$~rad) queda dentro de la región, a $0.028$~rad del borde. \\[3pt]
Evidencia &
\texttt{e1\_6\_region.m}, \texttt{src/simular\_lote.m}; figuras \texttt{figuras/E1\_6\_cortes.png} y
\texttt{figuras/E1\_6\_razon.png}; datos en \texttt{resultados/E1\_6.mat}. \\[3pt]
Comprobación independiente &
(1) $\max|Kx|$ en $10^5$ puntos de la frontera: $9.996$~V ($\le10$).
(2) $\dot V<0$ en la frontera para el modelo lineal y para el no lineal (100\,\% de los puntos); $\dot V\le0$
en todo el espacio, con convergencia por LaSalle: $(A,Q_0^{1/2})$ observable (rango 4).
(3) 5000 condiciones dentro del elipsoide, planta no lineal con retención de orden cero: 5000
recuperadas, $\max|u_c|=9.815$~V, ninguna muestra saturada.
(4) El simulador por lotes coincide con \texttt{simular\_qube} en 60 de 60 condiciones (diferencia final
$\le1.7\times10^{-19}$).
(5) Radio sobre $\alpha$ confirmado con el simulador oficial del kit (5 semillas). \\[3pt]
Supuestos &
El certificado vale para el modelo lineal en tiempo continuo, con estado exacto y sin límite de
recorrido; que se cumpla en la planta no lineal y muestreada es evidencia numérica, no prueba. Es
conservador porque exige que el motor nunca sature y porque se limita a un conjunto de nivel de una
función de Lyapunov cuadrática. Las fracciones del Monte Carlo dependen de la caja declarada
($|\theta|\le\pi/2$, $|\alpha|\le0.8$~rad, velocidades $\le10$~rad/s, semilla 2026); la caja recorta la región en
$\theta$ (radio estimado 4.6~rad), por lo que la razón de volúmenes (5506) es una cota inferior. La bisección
supone que la región es estrellada desde el origen; cerca del borde, la cuantización hace el resultado
sensible (radio de recuperación oficial $0.457$~rad, mayor que el ideal). Horizonte de 5~s.
\end{tabularx}
\end{recuadro}
